\section*{अध्याय \ref{ch:b1:enzymes} --- एंजाइम और जैवरासायनिक उत्प्रेरण}

\begin{solution}{exo:b1:enzymes:1}
वह \omterm{prop:b1:enzymes:whatitdoes}{सक्रियण ऊर्जा} (\omterm{prop:b1:enzymes:whatitdoes}{संक्रमण अवस्था} की ऊँचाई) घटाता है, और इस तरह दर भी,
दोनों दिशाओं में; वह $\Delta G$, साम्य नियतांक और साम्य की स्थिति
अपरिवर्तित छोड़ देता है। परिच्छेदिका पर शिखर नीचा हो जाता है, दोनों सिरे
नहीं।
\end{solution}

\begin{solution}{exo:b1:enzymes:2}
$K_m$: आधे $V_{\max}$ पर क्रियाधार-सांद्रता (\unit{mol/L})।
$V_{\max}$: संतृप्त \omterm{def:b1:enzymes:enzyme}{क्रियाधार} पर दर (\unit{mol/s}, या \unit{\micro
mol/min})। $k_{\mathrm{cat}}$: संतृप्ति पर प्रति सेकंड प्रति \omterm{def:b1:enzymes:enzyme}{एंजाइम} बदले गए
क्रियाधार-अणु (\unit{s^{-1}})। $k_{\mathrm{cat}}/K_m$: दक्षता,
यानी नीचे \omterm{def:b1:enzymes:enzyme}{क्रियाधार} पर द्वितीय-कोटि दर-नियतांक
(\unit{L.mol^{-1}.s^{-1}})।
\end{solution}

\begin{solution}{exo:b1:enzymes:3}
कोई \omterm{def:b1:enzymes:inhibitors}{संदमक} नहीं: अंतःखंड 1, इसलिए $V_{\max} = 1$; $x$ अंतःखंड $-0.5$,
इसलिए $K_m = 2$। स्पर्धी: $V_{\max} = 1$, $x$ अंतःखंड $-0.25$, $K_m = 4$।
अस्पर्धी: अंतःखंड 2, $V_{\max} = 0.5$; $x$ अंतःखंड $-0.5$,
$K_m = 2$।
\end{solution}

\begin{solution}{exo:b1:enzymes:4}
\omterm{def:b1:proteins:allostery}{अन्यस्थानिक} (पुनर्भरण) नियमन, मिलीसेकंड; \omterm{prop:b1:enzymes:control}{सहसंयोजी रूपांतरण}
(फ़ॉस्फ़ोरिलन), सेकंड से मिनट; किसी ज़ाइमोजन का प्रोटीन-अपघटनी सक्रियण,
एक ही बार, माँग पर; जीन-अभिव्यक्ति तथा अपघटन से मात्रा का नियंत्रण,
मिनटों से घंटों तक।
\end{solution}

\begin{solution}{exo:b1:enzymes:5}
$v_0 = 50[S]/(0.1 + [S])$: 8.3, 25, 45.5 और \qty{49.5}{\micro mol/min}।
$0.9 = [S]/(0.1 + [S])$ से $[S] = 0.9\,K_m/0.1 = \qty{0.9}{mmol/L}$ मिलता है।
\end{solution}

\begin{solution}{exo:b1:enzymes:6}
$e^{20\,000/2580} = e^{7.75} = 2300$। $10^{10}$ के लिए: $\Delta E_a =
RT\ln 10^{10} = 2.58\times 23.0 = \qty{59}{kJ/mol}$।
\end{solution}

\begin{solution}{exo:b1:enzymes:7}
\qty{1}{mL} में प्रति मिनट $V_{\max} = \qty{0.12}{mol/L}$:
यानी उत्पाद का \qty{1.2e-4}{mol/min} $= \qty{2.0e-6}{mol/s}$; \omterm{def:b1:enzymes:enzyme}{एंजाइम}
\qty{2e-12}{mol}: $k_{\mathrm{cat}} = 2\times 10^{-6}/2\times 10^{-12} =
\qty{1e6}{s^{-1}}$। दक्षता $10^6/0.012 = \qty{8e7}{L.mol^{-1}.s^{-1}}$,
यानी विसरण-सीमा के दस के गुणक के भीतर: क्रियाधार-अणु से लगभग हर मुलाक़ात
फलदायी है।
\end{solution}

\begin{solution}{exo:b1:enzymes:8}
$1 + 2/K_i = 2$: $K_i = \qty{2}{mmol/L}$। \qty{90}{\%} के लिए:
$[S] = 9K_m^{\mathrm{app}}$: \omterm{def:b1:enzymes:inhibitors}{संदमक} के बिना \qty{9}{mmol/L}, उसके साथ
\qty{18}{mmol/L}।
\end{solution}

\begin{solution}{exo:b1:enzymes:9}
पहला प्रतिबद्ध चरण संदमित करने से समूचा पथ एक साथ रुक जाता है और कोई
मध्यवर्ती बरबाद नहीं होता, जो अन्यथा जमा हो जाते; और जो मायने रखता है वह
संकेत अंत-उत्पाद है --- उसकी प्रचुरता ही ठीक वह सूचना है कि उस पथ की अब
ज़रूरत नहीं --- जबकि किसी मध्यवर्ती का स्तर माँग के बारे में कुछ नहीं कहता।
\end{solution}

\begin{solution}{exo:b1:enzymes:10}
जिस \omterm{def:b1:cell-unit-of-life:cell}{कोशिका} में कोई प्रोटीन-अपघटनी बने उसी में सक्रिय होकर वह उसी \omterm{def:b1:cell-unit-of-life:cell}{कोशिका}
को पचा डालती; निष्क्रिय ज़ाइमोजन के रूप में बनी हुई वह तब तक हानिरहित
रहती है जब तक आँत की अंतःस्तर का \omterm{def:b1:enzymes:enzyme}{एंजाइम} एंटरोपेप्टिडेज़ ट्रिप्सिनोजन को
काटकर ग्रहणी में ट्रिप्सिन न बना दे, जहाँ पाचन चाहिए; फिर ट्रिप्सिन बाकी
ज़ाइमोजन सक्रिय करती है (एक अनुक्रमिका)। और चूँकि ट्रिप्सिन का एक अंश भी
उस अनुक्रमिका को समय से पहले शुरू कर सकता है, अग्न्याशय बीमे के तौर पर
अपनी कणिकाओं में एक विशिष्ट ट्रिप्सिन-संदमक बाँध देता है; इस व्यवस्था की
विफलता ही अग्न्याशयशोथ है।
\end{solution}

\begin{solution}{exo:b1:enzymes:11}
बिना: $3^3/(3^3 + 3^3) = 0.50$। साथ: $27/(166 + 27) = 0.14$: यानी
दर अपने मान की \qty{28}{\%} रह जाती है। $K_m$ 3 से 5.5 तक चढ़ने वाला कोई
माइकेलिस--मेंटेन \omterm{def:b1:enzymes:enzyme}{एंजाइम}: $3/6 = 0.50$ से $3/8.5 = 0.35$, यानी
\qty{70}{\%}। \omterm{def:b1:proteins:allostery}{सहकारिता} वक्र की उसी खिसकन को कार्यकारी बिंदु के पास दुगने
से अधिक असरदार बना देती है: यानी कोई मंदक नहीं, कोई
स्विच।
\end{solution}

\begin{solution}{exo:b1:enzymes:12}
कोई नंगा उत्प्रेरक (कोई प्लैटिनम पृष्ठ, कोई अम्ल) कई अभिक्रियाएँ
बिना भेद किए तेज़ करता है; कोई \omterm{def:b1:enzymes:enzyme}{एंजाइम} एक \omterm{def:b1:enzymes:enzyme}{क्रियाधार} पर एक ही अभिक्रिया
तेज़ करता है --- उसके \omterm{def:b1:enzymes:enzyme}{सक्रिय स्थल} की विशिष्टता ही पहले से इस बारे में
सूचना है कि \omterm{def:b1:cell-unit-of-life:cell}{कोशिका} क्या करवाना चाहती है। नियमन एक दूसरी परत जोड़ता है:
\omterm{def:b1:proteins:allostery}{अन्यस्थानिक} स्थल, फ़ॉस्फ़ोरिलन स्थल और ज़ाइमोजन-पेप्टाइड \omterm{def:b1:enzymes:enzyme}{एंजाइम} को
\omterm{def:b1:cell-unit-of-life:cell}{कोशिका} की अवस्था के जवाब में बताते हैं कि कब और कहाँ काम करना है। लागत एक
बड़ा \omterm{def:b1:proteins:peptide}{प्रोटीन} है (एक दर्जन के स्थल के लिए सैकड़ों अवशेष), सबसे अच्छे
अकार्बनिक उत्प्रेरकों से धीमा आवर्तन, और उसे बनाने, बदलने तथा नष्ट करने
का तंत्र; लाभ ऐसा रसायन है जो केवल वहीं और तभी चलता है जहाँ और जब वह काम
का है, और उपापचय यही है।
\end{solution}

\begin{solution}{pb:b1:enzymes:1}
\textbf{1.} $1/[S]$: 2, 1, 0.5, 0.2, 0.1, 0.05. $1/v_0$: 5.0, 3.0, 2.0,
1.4, 1.2, 1.1.
\textbf{2.} $1/[S]$ का हर चरण $1/v_0$ को अनुपात में बदलता है: ढाल
$(5.0 - 1.1)/(2 - 0.05) = 2.0$; अंतःखंड $1.0$।
\textbf{3.} $V_{\max} = 1/1.0 = \qty{1.0}{\micro mol/min}$; $K_m =
\text{ढाल}\times V_{\max} = \qty{2.0}{mmol/L}$।
\textbf{4.} $1.0\times 5/(2 + 5) = 0.714$: जैसा नापा गया।
\textbf{5.} $\qty{10e-9}{mol/L}\times\qty{1e-3}{L} = \qty{1e-11}{mol}$;
$V_{\max} = \qty{1e-6}{mol/min} = \qty{1.67e-8}{mol/s}$;
$k_{\mathrm{cat}} = 1.67\times 10^{-8}/10^{-11} = \qty{1670}{s^{-1}}$।
\textbf{6.} $1670/(2\times 10^{-3}) = \qty{8.3e5}{L.mol^{-1}.s^{-1}}$:
यानी विसरण-सीमा से हज़ार गुना नीचे; हज़ार में केवल एक टक्कर
फलदायी है।
\textbf{7.} $0.95 = [S]/(2 + [S])$: $[S] = 19\,K_m = \qty{38}{mmol/L}$।
\textbf{8.} $1/v_0$: 9.0, 5.0, 3.0, 1.8, 1.4, 1.2. ढाल $(9.0 -
1.2)/1.95 = 4.0$; अंतःखंड $1.0$।
\textbf{9.} $V_{\max}^{\mathrm{app}} = \qty{1.0}{\micro mol/min}$
(अपरिवर्तित); $K_m^{\mathrm{app}} = 4.0\times 1.0 = \qty{4.0}{mmol/L}$।
\textbf{10.} $K_m$ दुगना, $V_{\max}$ अपरिवर्तित: स्पर्धी।
\textbf{11.} $4 = 2(1 + 1/K_i)$: $K_i = \qty{1.0}{mmol/L}$।
\textbf{12.} \omterm{def:b1:enzymes:inhibitors}{संदमक} के साथ, $v = V_{\max}[S]/(K_m(1 + [I]/K_i) +
[S])$; दर आधी करने के लिए $K_m(1 + [I]/K_i) + [S] = 2(K_m +
[S])$ चाहिए, यानी $[I] = K_i(K_m + [S])/K_m$: \qty{2}{mmol/L} पर, $[I] =
1\times 4/2 = \qty{2}{mmol/L}$; \qty{20}{mmol/L} पर, $[I] = 22/2 =
\qty{11}{mmol/L}$। जितना अधिक \omterm{def:b1:enzymes:enzyme}{क्रियाधार}, उतना अधिक \omterm{def:b1:enzymes:inhibitors}{संदमक} लगता है।
\textbf{13.} \omterm{def:b1:enzymes:enzyme}{क्रियाधार} (या उसकी \omterm{prop:b1:enzymes:whatitdoes}{संक्रमण अवस्था}) जैसा कोई अणु
--- कोई एस्टर-सदृश जो जल-अपघटित न हो सके --- जो
\omterm{def:b1:enzymes:enzyme}{सक्रिय स्थल} में बँधता हो।
\textbf{14.} $1/v_0$: 10.0, 6.0, 4.0, 2.8, 2.4, 2.2; ढाल $4.0$,
अंतःखंड $2.0$: $V_{\max}^{\mathrm{app}} = \qty{0.5}{\micro mol/min}$,
$K_m^{\mathrm{app}} = 4.0\times 0.5 = \qty{2.0}{mmol/L}$।
\textbf{15.} $V_{\max}$ आधा, $K_m$ अपरिवर्तित: अस्पर्धी;
$2 = 1 + 1/K_i$: $K_i = \qty{1.0}{mmol/L}$।
\textbf{16.} J \omterm{def:b1:enzymes:enzyme}{सक्रिय स्थल} के अलावा किसी स्थल पर, मुक्त \omterm{def:b1:enzymes:enzyme}{एंजाइम} और
$ES$ दोनों पर, उतनी ही बंधुता से बँधता है: \omterm{def:b1:enzymes:enzyme}{क्रियाधार} उससे होड़ नहीं करता,
और किसी भी $[S]$ पर आधे एंजाइम-अणु निष्क्रिय रहते हैं।
\textbf{17.} J के साथ दर दुगनी हो जाती है (दुगने \omterm{def:b1:enzymes:enzyme}{एंजाइम} का आधा भी
सक्रिय ही है); I के साथ भी दुगनी हो जाती है (दर हर $[S]$ पर
\omterm{def:b1:enzymes:enzyme}{एंजाइम} के अनुपात में है): \omterm{def:b1:enzymes:enzyme}{एंजाइम} दुगना करने से दोनों \omterm{def:b1:enzymes:inhibitors}{संदमक} कभी अलग नहीं
पहचाने जाते।
\textbf{18.} \qty{25}{\celsius} पर दर: $V_{\max}[S]/(K_m + [S])$ जिसमें
$V_{\max} = k_{\mathrm{cat}}[E]$: प्रति लीटर, $1670\times 10^{-8} =
\qty{1.67e-5}{mol/L/s}$; $\times 0.5/2.5 = \qty{3.3e-6}{mol/L/s}$;
\qty{37}{\celsius} पर, $\times 2^{1.2} = 2.3$: \qty{7.7e-6}{mol/L/s};
\qty{1e-12}{L} में: \qty{7.7e-18}{mol/s} $= \num{4.6e6}$ अणु प्रति
सेकंड।
\textbf{19.} किसी चढ़ाई की ज़रूरत नहीं: दर (\num{4.6e6}) माँग
(\num{3e6}) से अधिक है; और यदि उसे चढ़ना ही पड़ता, तो \omterm{def:b1:cell-unit-of-life:cell}{कोशिका} \omterm{def:b1:enzymes:enzyme}{क्रियाधार}
चढ़ा सकती थी, या \omterm{def:b1:enzymes:enzyme}{एंजाइम} को फ़ॉस्फ़ोरिलित करके उसका $K_m$ गिरा सकती थी, या
अधिक \omterm{def:b1:enzymes:enzyme}{एंजाइम} बना सकती थी।
\textbf{20.} $K_m^{\mathrm{app}} = 2(1 + 0.4/0.2) = \qty{6}{mmol/L}$: दर
$\times 0.5/6.5$, न कि $0.5/2.5$: यानी पहले की \qty{38}{\%},
प्रति सेकंड \num{1.8e6} --- उत्पाद अपने ही संश्लेषण को माँग से नीचे गला
देता है, इसलिए वह खर्च होता है और उसका स्तर गिरता है, जो
\omterm{def:b1:enzymes:enzyme}{एंजाइम} को छोड़ देता है: यानी एक स्वयं-समायोजी आपूर्ति।
\textbf{21.} पहले: $0.5/2.5 = 0.20$ गुना $V_{\max}$; बाद में: $0.5/0.9 =
0.56$: यानी लगभग तीन गुना।
\textbf{22.} लगभग शून्य: pH 5 पर \omterm{def:b1:enzymes:enzyme}{सक्रिय स्थल} का हिस्टिडीन
प्रोटॉनित रहता है और क्षार का काम नहीं कर सकता, और स्थल के अम्लीय अवशेष
उदासीन हो जाते हैं; \omterm{def:b1:enzymes:enzyme}{एंजाइम} की आयनन-अवस्था, और शायद उसका वलन भी,
गलत हैं। वह \omterm{def:b1:cell-unit-of-life:cell}{कोशिकाद्रव्य} के लिए गढ़ा है, \omterm{def:b1:eukaryotic-cell:endomembrane}{लयनकाय} के लिए नहीं।
\textbf{23.} $K_m$ थोड़ा ही बदलता है (बँधना अधिकतर स्थल के बाकी हिस्से
से होता है); $k_{\mathrm{cat}}$ परिमाण की कई कोटियों से ढह जाता है, चूँकि
हिस्टिडीन ही वह उत्प्रेरक क्षार था जो जल या सेरीन को सक्रिय करता था।
\textbf{24.} $K_m$ के पास दर \omterm{def:b1:enzymes:enzyme}{क्रियाधार} के बदलावों पर जवाब देती है
(ढाल $V_{\max}/2K_m$ के पास); उससे बहुत ऊपर \omterm{def:b1:enzymes:enzyme}{एंजाइम} संतृप्त और
असंवेदी रहता है, और अतिरिक्त \omterm{def:b1:enzymes:enzyme}{क्रियाधार} एक \omterm{def:b1:membranes-transport:osmosis}{परासरणी} तथा रासायनिक
बोझ होता। $K_m$ के पास काम करने से पथ आपूर्ति से नियंत्रणीय
बना रहता है।
\textbf{25.} $K_m = \qty{2.0}{mmol/L}$, $V_{\max} = \qty{1.0}{\micro mol/min}$
(\qty{10}{nmol/L} \omterm{def:b1:enzymes:enzyme}{एंजाइम}), $k_{\mathrm{cat}} = \qty{1670}{s^{-1}}$; I
स्पर्धी, $K_i = \qty{1.0}{mmol/L}$; J अस्पर्धी, $K_i =
\qty{1.0}{mmol/L}$।
\end{solution}
